% Seeded from ../slides/lecture_32.tex; substantive frame content retained.
% Title and primary objective become handout metadata; presentation commands removed.
\documentclass[compact]{ne630handout}
\handoutsolutionsfalse
% Assigned reading: ../administrivia/schedule-source.psv.
\begin{document}
\Lesson{32}{One-Speed Diffusion in a Homogeneous Slab}{FNRP 6.3.}{Students will be able to develop solutions for the one-speed diffusion equation in slab geometry for single-region systems with constant $D$, $\Sigma_a$, and a uniform volumetric source $S(x)$}
\begin{handoutflow}


% Source: lecture_32.tex, frame at line 89.
\Sec{1}{Review: from Continuity to Diffusion}
Last time, we developed the neutron transport equation and integrated it
over all angles to produce the neutron continuity equation
\begin{align*}
  \frac{1}{v}\,\frac{\partial \phi}{\partial t}
  &+ \nabla\!\cdot\! \mathbf{J}(\mathbf{r},E,t)
   +\Sigma_t(\mathbf{r},E)\, \phi(\mathbf{r},E,t) = \overbrace{\textcolor{wcprimary}{Q(\mathbf{r}, E, t)}}^{\text{in-scatter + source}} \, ,
  \tag{4--79}
\end{align*}
where $\phi$ is the (scalar) flux, $\mathbf{J}$ is the current density, and
the {\it emission density} $Q$ represents all \textcolor{wcprimary}{gains}.
Closure of the continuity equation  comes from Fick's law,
\begin{equation*}
\tag{FNRP 6.11}
   \mathbf{J} = -D \nabla \phi \, ,
\end{equation*}
substitution of which into (4--79) produces the {\bf neutron diffusion equation}
\begin{align*}
  \frac{1}{v}\,\frac{\partial \phi}{\partial t}
  - \nabla\!\cdot\! D(\mathbf{r},E) \nabla \phi(\mathbf{r},E,t)
   +\Sigma_t(\mathbf{r},E)\, \phi(\mathbf{r},E,t)
  = Q(\mathbf{r}, E, t)\, .
\end{align*}

% Source: lecture_32.tex, frame at line 117.
\Sec{2}{Review: Steady-State, One-Speed Diffusion}
Now that we have space, it's time to say goodbye to $t$ and to
replace $E$ by $g$:
\begin{align*}
 - \nabla\!\cdot\! D_g(\mathbf{r}) \nabla \phi_g(\mathbf{r})
   + \Sigma_{tg}(\mathbf{r})\, \phi_g(\mathbf{r})
  = \sum_{g'=1}^{G} \Sigma_{sg\gets g'} \phi_{g'}(\mathbf{r}) +  S_g(\mathbf{r}) \, .
\end{align*}

{\bf Example}: For $G=2$ and no upscatter, show that
\begin{align*}
 - \nabla\!\cdot\! D_1(\mathbf{r}) \nabla \phi_1(\mathbf{r})
   + \Sigma_{r,1}(\mathbf{r})\, \phi_1(\mathbf{r})
  &=  S_1(\mathbf{r},t) \\
 - \nabla\!\cdot\! D_2(\mathbf{r}) \nabla \phi_2(\mathbf{r})
   + \Sigma_{a,2}(\mathbf{r})\, \phi_2(\mathbf{r})
  &=  \Sigma_{s,2\gets 1} \phi_{1}(\mathbf{r}) + S_2(\mathbf{r})
\end{align*}

What if $G=1$ and we introduce multiplication?

\begin{align*}
\tag{FNRP 6.12}
\boxed{ - \nabla\!\cdot\! D(\mathbf{r}) \nabla \phi(\mathbf{r})
   + \Sigma_{a}(\mathbf{r})\, \phi(\mathbf{r})
  =  S(\mathbf{r}) + \bar{\nu}\Sigma_{f}(\mathbf{r})\, \phi(\mathbf{r})
}
\end{align*}

% Source: lecture_32.tex, frame at line 153.
\Sec{3}{Applying $\nabla$ and $\nabla^2$}
Much like integration over $4\pi$ hides some detail, so, too, does $\nabla\!\cdot\! D(\mathbf{r}) \nabla \phi(\mathbf{r})$!

For Cartesian (or $xyz$) coordinates, the gradient and divergence operators act on scalar and vector functions according to
\begin{align*}
\nabla \phi(\mathbf{r})
= \mathbf{i} \,\frac{\partial \phi}{\partial x}
 + \mathbf{j}  \,\frac{\partial \phi}{\partial y}
 + \mathbf{k}  \,\frac{\partial \phi}{\partial z} \quad \text{and} \quad
\nabla\!\cdot\!\mathbf{J}(\mathbf{r})
= \frac{\partial J_x}{\partial x}
 + \frac{\partial J_y}{\partial y}
 + \frac{\partial J_z}{\partial z} \, ,
\end{align*}
so that
\begin{align*}
\nabla\!\cdot\!\big(D(\mathbf{r})\,\nabla \phi(\mathbf{r})\big)
&= \frac{\partial}{\partial x}\!\left(D\,\frac{\partial \phi}{\partial x}\right)
 + \frac{\partial}{\partial y}\!\left(D\,\frac{\partial \phi}{\partial y}\right)
 + \frac{\partial}{\partial z}\!\left(D\,\frac{\partial \phi}{\partial z}\right) \\[4pt]
&= D(\mathbf{r})\,\nabla^2 \phi(\mathbf{r})
  + \big(\nabla D(\mathbf{r})\big)\!\cdot\!\big(\nabla \phi(\mathbf{r})\big).
\end{align*}

% Source: lecture_32.tex, frame at line 188.
\Sec{4}{The Homogeneous Slab}
For homogeneous regions (i.e., $\Sigma_a(\mathbf{r}) = \Sigma_a$ and $D(\mathbf{r}) = D$),
the one-speed diffusion equation can be written as
\begin{align*}
  -D \left (\frac{\partial^2 \phi}{\partial x^2} +
            \frac{\partial^2 \phi}{\partial y^2} +
            \frac{\partial^2 \phi}{\partial z^2} \right )
  + \Sigma_{a}\phi(\mathbf{r})
  = \bar{\nu}\Sigma_{f} \phi(\mathbf{r})  +  S(\mathbf{r}) \, .
\end{align*}

Now, suppose this region is very (infinitely) large in the $y$ and $z$ directions but finite along the $x$.  If $S(\mathbf{r}) = S(x)$, then everything in this {\it homogeneous slab} is uniform along the $y$ and $z$ directions, and, hence, we should have $\phi(\mathbf{r}) = \phi(x)$.  Thus, the one-speed diffusion equation in slab geometry is
\begin{align*}
  -D \frac{d^2 \phi}{d x^2}
  + \Sigma_{a}\phi(x)
  = \bar{\nu}\Sigma_{f} \phi(x)  +  S(x) \, .
\end{align*}

This is a second-order, inhomogeneous, ordinary differential equation with constant coefficients.

% Source: lecture_32.tex, frame at line 215.
\Sec{5}{The Diffusion Coefficient $D$}
The {\bf diffusion coefficient} is defined\footnote{through the formal derivation of the diffusion equation from the transport equation} as
\begin{equation*}
 D \;=\; \frac{1}{3\,\Sigma_{\mathrm{tr}}} \, ,
\end{equation*}
where $\Sigma_{\mathrm{tr}} \equiv \Sigma_t - \bar{\mu}\,\Sigma_s$ is the {\bf transport cross section}, and
$\bar{\mu}$ is the average cosine of the scattering angle in the {\it laboratory} system.

For isotropic elastic scattering in the
{\it center-of-mass system} and a single
nuclide of mass number $A$,
\begin{equation*}
\tag{FNRP 2.X}
 \bar{\mu} = \frac{2}{3A} \, .
\end{equation*}
When neither $\Sigma_{\mathrm{tr}}$ nor $\bar{\mu}$ is specified, we set $\Sigma_{\mathrm{tr}} \approx \Sigma_{t}$.

{\bf Ponderable}: What does $\Sigma_{\mathrm{tr}} \approx \Sigma_{t}$ imply about scattering?

% Source: lecture_32.tex, frame at line 243.
\Sec{6}{The Diffusion Length $L$ and A First Example}
To simplify things a bit, let's define the {\bf diffusion length}
\begin{equation*}
 L = \sqrt{\frac{D}{\Sigma_a}} \, ,
\end{equation*}
so that the diffusion equation sans multiplication becomes
\begin{align*}
   \frac{d^2 \phi}{d x^2}
  - \frac{1}{L^2}\phi(x)
  = -\frac{1}{D} S(x) \, .
\end{align*}

{\bf Example}. Consider a slab of thickness $a$ with known $\Sigma_a$ and $D$. If there is a uniform source $S_0$ throughout, determine $\phi(x)$ in $0\le x\le a$ subject to
the boundary conditions\footnote{We'll see other boundary conditions next time, but the zero-flux conditions are an approximate way to model a slab surrounded by ``nothing'' (i.e., a vacuum}
\begin{equation*}
 \phi(0) = 0  \qquad \text{and} \qquad \phi(a) = 0 \, .
\end{equation*}

% Source: lecture_32.tex, frame at line 265.
\Sec{7}{A First Example: Solution}
\textbf{Step 1 (Write the Equations)}. We have (with $L=\sqrt{D/\Sigma_a}$)
\begin{equation*}
\tag{ODE and BCs}
 \frac{d^2\phi}{dx^2} - \frac{1}{L^2}\,\phi = -\frac{S_0}{D}  \quad \text{subject to} \quad  \phi(0) = 0  \,\, \text{and} \,\, \phi(a) = 0 \, .
\end{equation*}
\textbf{Step 2 (Homogeneous Solution)}. Try $\phi_h(x)=e^{\lambda x}$ to obtain $\lambda=\pm L^{-1}$, and, hence,
\begin{equation*}
 \phi_h(x)=C_1 e^{x/L}+C_2 e^{-x/L}.
\end{equation*}
\textbf{Step 3 (Particular Solution)}. Since $S_0$ is constant, take $\phi_p(x)=C_3$; substitution into ODE gives $C_3=L^2 S_0 /D = S_0/\Sigma_a$.\\[4pt]
\textbf{Step 4 (Apply BCs)}. Plugging the general solution $\phi_h+\phi_p$ into the BCs leads to
\begin{align*}
\tag{left BC}
 \phi(0) &=\frac{S_0}{\Sigma_a}+C_1 e^{0/L}+C_2 e^{-0/L} = 0\\
\tag{right BC}
 \phi(a) &=\frac{S_0}{\Sigma_a}+C_1 e^{a/L}+C_2 e^{-a/L} = 0 \, ,
\end{align*}
which represent two equations for the two unknowns, $C_1$ and $C_2$.  {\bf Find them!}

% Source: lecture_32.tex, frame at line 291.
\Sec{8}{A First Example: Verification and Inspection}
We can write the solution compactly\footnote{We can define any two independent, linear combinations
of $e^{x/L}$ and $e^{-x/L}$, say $f_{\pm}(x) =  (e^{x/L} \pm e^{-x/L})/2$
instead define $\phi_h(x) = \tilde{C}_1 f_+(x) + \tilde{C}_2 f_-(x)$.  For many slab problems,
we'll find that these functions $f_+(x) = \cosh(x/L)$ and $f_-(x) =\sinh(x/L)$ can produce
simpler solutions.
}
\begin{equation*}
 \tag{***}
 \phi(x)=\frac{S_0}{\Sigma_a}\left[1-\frac{\cosh\!\big(\tfrac{x-a/2}{L}\big)}{\cosh\!\big(\tfrac{a}{2L}\big)}\right] \, .
 \label{eq:ex1sol}
\end{equation*}
To verify this is correct, we should
\begin{enumerate}
 \item Plug (***) into the ODE.
 \item Plug (***) into the BCs.
 \item Plot $\phi(x)$ to inspect the shape and boundary values.
\end{enumerate}
When we aren't given numerical values for $S_0$, $\Sigma_a$, and $D$, we can simply assume reasonable values.

{\bf Example}: Plot $\phi(x)$ with $\Sigma_a=0.1$ \si{\per\centi\meter}, $D=1$ \si{\centi\meter}, $S_0=1$ \si{\per\centi\meter\per\second} for (a) $a=10$ \si{\centi\meter} and (b) $a=100$ \si{\centi\meter}.  What's the maximum possible flux for any $a$?

\end{handoutflow}
\end{document}
